\documentclass[12pt, twoside]{article}
\input{../preamble}

\fancyhead[LE]{\thepage}
\fancyhead[RO]{Name: \hspace{4cm} \,\\ \hspace{2.9cm} \,}
\fancyhead[LO]{La Scuola d'Italia / Dr.\ Huson / IB Math \\* 25 March 2026}

\begin{document}
\subsubsection*{5.6 Classwork: ``Show That'' Problems (no calculator)}

IB command term: \textbf{show that} --- obtain the required result without the formality of proof.\\[0.25cm] ``Show that'' questions are
\emph{not} verify; you may not simply substitute values and show that the equation holds. \\[0.25cm]
Instead, you must derive the given result with all of the algebraic steps, including the final, given result. \vspace{1cm}

\begin{enumerate}[itemsep=0.2cm]

%--- Problem 1: Kinematics integration (2021 Nov P1 Q7) ---
\item \textbf{[Maximum mark: 16]}

A particle $P$ moves along the $x$-axis. The velocity of $P$ is $v$\;m\,s$^{-1}$
at time $t$ seconds, where
\[v(t) = 4 + 4t - 3t^2 \quad \text{for } 0 \leq t \leq 3.\]
When $t = 0$, $P$ is at the origin $O$.

\begin{enumerate}[itemsep=0.15cm]
  \item (i) \; Find the value of $t$ when $P$ reaches its maximum velocity.

        (ii) Show that the distance of $P$ from $O$ at this time is
             $\dfrac{88}{27}$ metres. \hfill [7]
  \item Sketch a graph of $v$ against $t$, clearly showing any points of
        intersection with the axes. \hfill [4]
  \item Find the total distance travelled by $P$. \hfill [5]
\end{enumerate}

\vfill
\hfill \emph{IB AA SL 2021 Nov P1 Q7}

%--- Problem 2: Quotient rule (2021 May P1 TZ1 Q8) ---
\item \textbf{[Maximum mark: 13]}

Let $\displaystyle y = \frac{\ln x}{x^4}$ for $x > 0$.

\begin{enumerate}[itemsep=0.15cm]
  \item Show that $\displaystyle \frac{dy}{dx} = \frac{1 - 4\ln x}{x^5}$. \hfill [3]
  \item The graph of $f$ has a horizontal tangent at point $P$.
        Find the coordinates of $P$. \hfill [5]
  \item Given that $\displaystyle f''(x) = \frac{20\ln x - 9}{x^6}$,
        show that $P$ is a local maximum point. \hfill [3]
  \item Solve $f(x) > 0$ for $x > 0$. \hfill [2]
\end{enumerate}

\vfill
\hfill \emph{Adapted from IB AA SL 2021 May P1 TZ1 Q8}


\end{enumerate}
\end{document}
